\section{Duality for compactly supported cohomology}
\label{section-duality-compactly-supported}

\noindent
Let $k$ be a field. Let $U$ be a separated scheme of finite type over $k$.
Let $K$ be an object of $D^b_{\textit{Coh}}(\mathcal{O}_U)$. Let us define
the compactly supported cohomology $H^i_c(U, K)$ of $K$ as follows.
Choose an open immersion $j : U \to X$ into a scheme proper over $k$
and a Deligne system $(K_n)$ for $j : U \to X$ whose restriction
to $U$ is constant with value $K$. Then we set
$$
H^i_c(U, K) = \lim H^i(X, K_n)
$$
We view this as a topological $k$-vector space using the limit topology
(see More on Algebra, Section \ref{more-algebra-section-topological-ring}).
There are several points to make here.

\medskip\noindent
First, this definition is independent of the choice of $X$ and $(K_n)$.
Namely, if $p : U \to \Spec(k)$ denotes the structure morphism, then
we already know that $Rp_!K = (R\Gamma(X, K_n))$
is well defined up to pro-isomorphism in $D(k)$ hence so is the limit
defining $H^i_c(U, K)$.

\medskip\noindent
Second, it may seem more natural to use the expression
$$
H^i(R\lim R\Gamma(X, K_n)) = R\Gamma(X, R\lim K_n)
$$
but this would give the same answer: since the $k$-vector spaces
$H^j(X, K_n)$ are finite dimensional, these inverse systems satisfy
Mittag-Leffler and hence $R^1\lim$ terms of Cohomology, Lemma
\ref{cohomology-lemma-RGamma-commutes-with-Rlim} vanish.

\medskip\noindent
If $U' \subset U$ is an open subscheme, then there is a canonical map
$$
H^i_c(U', K|_{U'}) \longrightarrow H^i_c(U, K)
$$
functorial for $K$ in $D^b_{\textit{Coh}}(\mathcal{O}_U)$.
See for example Remark \ref{remark-covariance-open-lower-shriek}.
In fact, using Remark \ref{remark-covariance-etale-lower-shriek}
we see that more generally such a map exists for an \'etale morphism
$U' \to U$ of separated schemes of finite type over $k$.

\medskip\noindent
If $V$ is a $k$-vector space then we put a topology on $\Hom_k(V, k)$
as follows: write $V = \bigcup V_i$ as the filtered union of its finite
dimensional $k$-subvector spaces and use the limit topology on
$\Hom_k(V, k) = \lim \Hom_k(V_i, k)$. If $\dim_k V < \infty$ then
the topology on $\Hom_k(V, k)$ is discrete. More generally, if
$V = \colim_n V_n$ is written as a directed colimit of finite dimensional
vector spaces, then $\Hom_k(V, k) = \lim \Hom_k(V_n, k)$ as topological
vector spaces.

